% Part:sets-functions-relations
% Chapter: size-of-sets
% Section: pairing

\documentclass[../../../include/open-logic-section]{subfiles}

\begin{document}

\olfileid{sfr}{siz}{pai}
\olsection{జతీకరణ ప్రమేయాలు, సంకేత సంఖ్యలు}

\begin{explain}
కాంటర్ జిగ్‌జాగ్ పద్ధతి $\Nat^n$ లెక్కించదగినదని కంటికి కనిపించేలా
చేస్తుంది. అయితే $\Nat^2$ను చూపే మన పట్టికపైనే దృష్టి పెడదాం.
పట్టికలో జిగ్‌జాగ్ గీత వెంబడి స్థానాలను లెక్కిస్తే,
$\tuple{1,2}$కు $7$ జతచేయబడిందని తనిఖీ చేయవచ్చు. కానీ దీన్ని మరింత
నేరుగా గణించగలిగితే బాగుంటుంది. అంటే, జిగ్‌జాగ్ లెక్కింపు యొక్క
\emph{విలోమం} మన దగ్గర ఉంటే బాగుంటుంది: $g\colon \Nat^2 \to \Nat$,
మరియు
\[
g(\tuple{0,0}) = 0, \;
g(\tuple{0,1}) = 1, \;
g(\tuple{1,0}) = 2, \; \dots,
g(\tuple{1,2}) = 7, \; \dots
\]
అయ్యేలా. అప్పుడు మన లెక్కింపులో $\tuple{n, m}$ కచ్చితంగా ఎక్కడ
వస్తుందో గణించగలం.

వాస్తవానికి రెండు విషయాలను గమనించి $g$ను నేరుగా నిర్వచించవచ్చు.
మొదటిది: $n$వ అడ్డువరుస, $m$వ నిలువువరుసలో విలువ $v$ ఉంటే,
$(n+1)$వ అడ్డువరుస, $(m-1)$వ నిలువువరుసలో విలువ $v + 1$ ఉంటుంది.
రెండవది: మన లెక్కింపు మొదటి అడ్డువరుసలో $0$, $1$, $3$, $6$ మొదలైన
త్రిభుజ సంఖ్యలు ఉంటాయి. $k$వ త్రిభుజ సంఖ్య అనేది kకు మించని
సహజ సంఖ్యల మొత్తం; దాన్ని $k(k+1)/2$గా గణించవచ్చు. ఈ రెండు
గమనింపులను కలిపి, ఈ ప్రమేయాన్ని పరిశీలించండి:
\[
  g(n,m) = \frac{(n+m+1)(n+m)}{2} + n
\]
$g(\tuple{n, m})$ అని రాయడం కంటే కంటికి సులభంగా ఉండటంతో తరచుగా
$g(n, m)$ అని మాత్రమే రాస్తాం. ముందుగా $(n+m)^\text{th}$ త్రిభుజ
సంఖ్యను కనుక్కొని, దానికి $n$ కలపాలని ఇది చెబుతుంది. ఇది మనకు
కావలసిన విధంగానే పట్టికను నింపుతుంది. ప్రత్యేకించి, జత
$\tuple{1, 2}$ను $\frac{4 \times 3}{2} + 1 = 7$కు పంపుతుంది.
\sourcecorrection{OLTESIZ-002}{ఆంగ్ల మూలం ``$k$ కంటే తక్కువ సహజ
సంఖ్యల మొత్తం'' అని చెప్పి $k(k+1)/2$ సూత్రాన్ని ఇస్తుంది. ఆ సూత్రం
$0+1+\cdots+k$కు సరిపోతుంది కనుక, ఇక్కడ ``$k$కు మించని'' అని
సరిచేశాం; జతీకరణ సూత్రాన్ని మార్చలేదు.}

ఈ ప్రమేయం $g$, జతల సమితి లెక్కింపు యొక్క \emph{విలోమం}. ఇటువంటి
ప్రమేయాలను \emph{జతీకరణ ప్రమేయాలు} అంటారు.
\end{explain}

\begin{defn}[జతీకరణ ప్రమేయం]
$f\colon A \times B \to \Nat$ ఒకటి-ఒకటి అయితే, $f$ను అంకగణిత
\emph{జతీకరణ ప్రమేయం} అంటాం. $f$, $A \times B$ను
\emph{సంకేతీకరిస్తుంది} అనీ, $f(x,y)$ను $\tuple{x,y}$ యొక్క
\emph{సంకేత సంఖ్య} అనీ కూడా అంటాం.
\end{defn}

\begin{explain}
ఉదాహరణకు, సహజ సంఖ్యల జతలను సంకేతీకరించడానికి జతీకరణ ప్రమేయాలను
వాడవచ్చు. మరోలా చెప్పాలంటే, ప్రతి మూలకాల \emph{జత}ను ఒకే
\emph{ఒక్క} సంఖ్యతో సూచించవచ్చు. జతీకరణ ప్రమేయపు విలోమాన్ని వాడి
ఆ సంఖ్యను \emph{విసంకేతీకరించి}, అది ఏ జతను సూచిస్తుందో తెలుసుకోవచ్చు.
\end{explain}

\begin{prob}
ఋణేతర కరణీయ సంఖ్యలన్నిటి సమితికి ఒక లెక్కింపును ఇవ్వండి.
\end{prob}

\begin{prob}
$\Rat$ \tetoken{లెక్కించదగిన}{!!{enumerable}} అని చూపండి. ఏ కరణీయ సంఖ్యనైనా $z \in \Int$,
$m \in \Nat^+$ అయ్యే $z/m$ భిన్నంగా రాయవచ్చని గుర్తుంచుకోండి.
\end{prob}

\begin{prob}
$\Bin^*$కు ఒక లెక్కింపును నిర్వచించండి.
\end{prob}

\begin{prob}
మీ పరిచయ తర్కశాస్త్ర పాఠ్యక్రమం నుంచి, సాధ్యమైన ప్రతి సత్యపట్టిక ఒక
సత్య ప్రమేయాన్ని వ్యక్తం చేస్తుందని గుర్తుంచుకోండి. అంటే, ఏదో ఒక
$k$కు సత్య ప్రమేయాలన్నీ $\Bin^k \to \Bin$ ప్రమేయాలు. సత్య
ప్రమేయాలన్నిటి సమితి లెక్కించదగినదని నిరూపించండి.
\end{prob}

\begin{prob}
ఏదైనా అనంత \tetoken{లెక్కించదగిన}{!!{enumerable}} సమితి యొక్క పరిమిత ఉపసమితులన్నిటి సమితి
\tetoken{లెక్కించదగిన}{!!{enumerable}} అని చూపండి.
\end{prob}

\begin{prob}
$\Nat$ యొక్క ఒక ఉపసమితి, $\Nat$లో తన పూరకం పరిమిత సమితి అయితేనే
దాన్ని \emph{పరిమిత-పూరక} సమితి అంటాం; అంటే, $A \subseteq \Nat$
పరిమిత-పూరక సమితి అవడం అంటే $\Nat\setminus A$ పరిమితంగా ఉండటం.
\sourcecorrection{OLSIZ-002}{ఆంగ్ల మూలంలోని అసంపూర్ణమైన ``complement
of a finite set $\Nat$'' వాక్యాన్ని, వెంటనే వచ్చే సమానార్థక నిబంధన
ప్రకారం ``$\Nat$లోని పరిమిత ఉపసమితికి $\Nat$లో పూరకం''గా స్పష్టం చేశాం.}
$\Nat$ యొక్క పరిమిత, పరిమిత-పూరక ఉపసమితులే కచ్చితంగా \tetoken{మూలకాలుగా}{!!{element}s}
గల సమితి $I$ అనుకుందాం. $I$ \tetoken{లెక్కించదగిన}{!!{enumerable}} అని చూపండి.
\end{prob}

\begin{prob}
\tetoken{లెక్కించదగిన}{!!{enumerable}} సమితుల \tetoken{లెక్కించదగిన}{!!{enumerable}} సమ్మేళనం \tetoken{లెక్కించదగిన}{!!{enumerable}} అని
చూపండి. అంటే $A_1$, $A_2$, \dots{} సమితులు, ఒక్కొక్క $A_i$
\tetoken{లెక్కించదగిన}{!!{enumerable}} అయినప్పుడల్లా, వాటన్నిటి సమ్మేళనం
$\bigcup_{i=1}^\infty A_i$ కూడా \tetoken{లెక్కించదగిన}{!!{enumerable}} అని చూపండి.
[గమనిక: ఇది కష్టం!]
\end{prob}

\begin{prob}
$f \colon A \times B \to \Nat$ ఏదైనా జతీకరణ ప్రమేయం అనుకుందాం.
$f^{-1}\colon \ran{f}\to A\times B$ ద్విగుణ ప్రమేయం అని చూపి,
$A\times B$ లెక్కించదగినదని తేల్చండి. అదనంగా $f$ అనేది $\Nat$పైకి
సంగ్రస్తమైతే, $f^{-1}\colon\Nat\to A\times B$ స్వయంగా ఒక లెక్కింపు
అని చూపండి.
\sourcecorrection{OLTESIZ-003}{ఆంగ్ల మూలం, ఒకటి-ఒకటి మాత్రమే కావలసిన
ఏ జతీకరణ ప్రమేయపు విలోమమైనా $A\times B$ లెక్కింపు అంటుంది. సాధారణంగా
విలోమపు ప్రవేశం $\ran{f}$ మాత్రమే; $\Nat$ అంతా కావాలంటే $f$ సంగ్రస్తం
కూడా కావాలి. ఇక్కడ రెండు సరైన నిర్ధారణలను వేరు చేశాం.}
\end{prob}

\begin{prob}
$\Nat^3$ను సంకేతీకరించే ఒక ప్రమేయాన్ని పేర్కొనండి.
\end{prob}

\end{document}
